\documentclass[11pt]{article}
\usepackage{icmlko}
\begin{document}

\title{기준선에 손잡이 스물넷을 더 줬더니 우리 배수가 반이 됐다 --- 오늘 아침에 낸 숫자를 오늘 저녁에 깎는다}
\author{월드모델 실험실}
\date{노트 584 · 2026년 8월 5일}
\maketitle

\begin{abstract}
노트 $580$ 이 ``통합 $\div$ 전용 $=$ 유보 가중 \textbf{$1.67$ 배}''를 내고
\textbf{덱과 논문 $449$ 에 실었다.} 그 노트가 스스로 한계를 적었다 ---
``전용 손잡이 격자가 넷뿐이다. 더 넓히면 전용이 오를 수 있다.''
\textbf{밖으로 나간 숫자이므로 그 한계를 먼저 닫았다.}

후보를 넷 $\to$ \textbf{스물여덟}(부스팅 $24$ $\cdot$ 랜덤포레스트 $2$
$\cdot$ 선형 $2$)로 넓혔다.

\begin{center}
\begin{tabular}{lrrrl}
\hline
 & 학습 & 후보 $4$ & \textbf{후보 $28$} & 고른 것 \\
\hline
시장 팝업 & $123$ & $\mathbf{13.84}$ & $\mathbf{1.07}$ & \textbf{랜덤포레스트} \\
아이돌 & $122$ & $1.36$ & $\mathbf{0.92}$ & \textbf{선형} \\
도서 & $80$ & $3.29$ & $3.80$ & 랜덤포레스트 \\
웹툰 & $2{,}938$ & $1.09$ & $1.22$ & 부스팅 \\
\hline
\textbf{유보 가중} & & $\mathbf{1.67}$ & $\mathbf{1.23}$ & \\
\hline
\end{tabular}
\end{center}

\textbf{비를 정한 것은 우리 모형이 아니라 기준선에 준 손잡이였다.}
사전 등록 셋 중 \textbf{둘이 깨졌다.}
\end{abstract}

\begin{figure}[t]
\centering
\includegraphics[width=\textwidth]{figs/f1.pdf}
\caption{\textbf{가} --- 나무 하나, 직선 하나가 뒤집는다.
\textbf{나} --- 기울기가 아니라 문턱이었다.
\textbf{다} --- 오늘 낸 숫자를 오늘 깎았다.}
\end{figure}

\section{사전 등록 판정}

\begin{center}
\begin{tabular}{lll}
\hline
 & 결과 & \\
\hline
① 가중 평균 $\geq 1.4$ 배 & $1.67 \to \mathbf{1.23}$ & \textbf{깨짐} \\
② 작을수록 크게 이긴다 & $\rho\,{-}0.70 \to \mathbf{-0.03}$ & \textbf{깨짐} \\
③ 지는 곳 $\leq 3$ & 둘(만화 $0.99$ $\cdot$ 아이돌 $0.92$) & 통과 \\
\hline
\end{tabular}
\end{center}

①이 깨지면 \textbf{덱을 다시 고친다}고 미리 적어 뒀다. 고쳤다.

\section{무엇이 비를 정했나}

큰 도메인 여덟은 \emph{전부} 부스팅을 고르고 비가 거의 안 움직였다
($\pm 0.03$). 뒤집힌 둘은 \textbf{다른 종류의 모형}을 골랐다.

\begin{itemize}
\item \textbf{시장 팝업} --- 랜덤포레스트가 $+0.0113 \to \mathbf{+0.1463}$.
비 $13.84 \to 1.07$. 옛 수치는 \textbf{$0$ 근처인 분모}가 만든 것이었다.
\item \textbf{아이돌} --- \textbf{선형 회귀}가 $+0.3458 \to \mathbf{+0.5090}$
으로 \textbf{통합을 이겼다.} 유보 $51$ 행이라 씨앗 잡음과 구별은 안 된다.
\end{itemize}

\subsection*{넓은 격자가 늘 좋은 것도 아니다}

웹툰 전용은 $+0.3283 \to \mathbf{+0.2929}$ 로 \emph{내려갔고}(비
$1.09 \to 1.22$), 게임 $-0.004$ $\cdot$ 도서 $-0.017$ $\cdot$ 모바일
$-0.015$ 도 내려갔다.

\begin{quote}
\textbf{교차검증은 오르는데 유보가 내린다.} 작은 시장에서는 \emph{손잡이
고르기 자체가} 과적합한다 --- 도서 CV $+0.092 \to +0.229$ 인데 유보
$+0.1227 \to +0.1061$. \textbf{전용 모형은 손잡이를 고를 검증 표본조차
없다.}
\end{quote}

\section{기울기가 아니라 문턱이었다}

덱 슬라이드 $6$ 의 ``읽는 법''이 ``\textbf{작은 시장일수록 더 크게
이깁니다}''였다. $\rho = -0.03$ --- \textbf{관계가 없다.} 학습 $122$ 행
위로는 $0.92 \sim 1.22$ 로 평평하다.

그 인상은 \textbf{시장 팝업 $13.84$ 와 도서 $3.29$ 두 칸}이 만들고
있었고, \textbf{하나가 통째로 무너졌다.}

\begin{quote}
남는 것은 \textbf{문턱}에 가깝다 --- 도서 $80$ 행 $3.80$ 배, 팝업 $24$ 행은
\emph{전용 모형이 성립하지 않는다}, $122$ 행 위는 평평. \textbf{다만 문턱
아래 시장이 둘뿐이다.} 덱에 그렇게 적었다.
\end{quote}

\section{격자와 무관하게 살아남는 것}

\begin{quote}
\textbf{팝업은 전용 모형을 만들 수가 없다.} 학습 $24$ 행이라 $5$겹이 안
선다. 후보를 몇 개 주든 마찬가지다. \textbf{통합은 그래도 $+0.3049$ 를
낸다.} 이것이 제일 센 문장이고 \emph{유일하게} 기준선 노력과 무관하다.
\end{quote}

그리고 통합이 $\mathbf{8/10}$ 에서 이긴다 $\cdot$ 가중 $\mathbf{1.23}$ 배
$\cdot$ 도서 $\mathbf{3.80}$ 배.

\section*{같은 함정을 네 번째로 만났다}

\begin{center}
\begin{tabular}{lll}
\hline
 & 잰 것 & 읽은 것 \\
\hline
노트 $548$ & 검색 축 전체의 값(상한) & 펀딩 수집의 기대값 \\
노트 $560$ & 전역 변경의 도메인별 차 & 그 도메인만 바꾼 이득 \\
노트 $568$ & 단일 특성의 유보 $\rho$ & 모형이 벌 양 \\
\textbf{노트 $584$} & \textbf{인색한 기준선 대비 비} & \textbf{모형의 값} \\
\hline
\end{tabular}
\end{center}

\textbf{앞의 셋과 달리 이번 것은 밖으로 나간 문서에 실렸다.} 그래서
조항을 사람이 기억하게 두지 않고 코드로 옮겼다 ---
\texttt{sideaudit.baseline\_ratio(비, 이름, 후보수, 유보행수)} 가 후보
여덟 아래 $\cdot$ 비 $3$ 초과 $\cdot$ 유보 $200$ 행 미만에 각각 다른
경고를 낸다.

\section*{조항}

\begin{quote}
\textbf{기준선 대비 비를 적을 때는 기준선에 준 노력을 같이 적는다.}
후보 넷과 후보 스물여덟은 다른 실험이고, 그 차이가 $13.84$ 와 $1.07$ 을
가른다. \textbf{``우리가 $N$ 배''는 ``우리가 무엇보다 $N$ 배''를 적기
전에는 문장이 아니다.}
\end{quote}

\section*{남은 것}

\begin{itemize}
\item \textbf{아이돌 $0.92$} --- 유보 $51$ 행이라 씨앗 잡음과 구별이 안
된다. 씨앗을 늘려 다시 재야 한다.
\item \textbf{전용에 안 준 것 셋} --- 행 무게(\texttt{TRAINW} $\cdot$
\texttt{TRAINR})와 \texttt{SEASON} $\cdot$ \texttt{SPEC} 칸. 통합에만 있는
것이 그 셋뿐이므로 그것까지 주면 \textbf{$1.23$ 도 더 내려갈 수 있다.}
\item \textbf{문턱 가설} --- 표본이 둘뿐(도서 $80$ $\cdot$ 팝업 $24$).
작은 시장을 더 만들어야 판정된다.
\end{itemize}

\end{document}
